% Pursuing Stacks (cotes 134-2 to 134-8): archive cote 134-6, site page 30 (PDF page 31), typescript page 348.
% typescript. Draft compared: ps5.tex:2810-2882.
% AI-transcribed by Claude Opus 5.5 (claude-opus-5-5, adaptive thinking, effort high), Message Batches API, batch ps3 (msgbatch_01XdFmuVVBAcgc7EFxdtCqsj), 2026-11-12; not yet checked by eye.
\scanpage{134-6/30}
\mspage{348}

b)  Assume the contractibility structure $M_c$  non trivial,  i.e.
there exists an interval
\[ \underline{I} = (I,\hand{$\delta$}_0,\hand{$\delta$}_1) \quad \text{with} \quad I \in M_c \ , \ \hand{$\delta$}_0 \hand{$\neq$} \hand{$\delta$}_1 \ . \]
Then for any small category A and any functor
\[ i\colon A \longrightarrow M \]
factoring through $M_c$ ,  the interval $i^*(\underline{I})$ in $A^\wedge$ is separating.  Hence
if  (for a given basic localizer $\underline{W}$) i is totally $\underline{W}$-aspheric (cf. theor\hand{em}
1 cor. 1  p.252),  hence $i^*(I)$ is totally aspheric in $A^\wedge$, \inserted{then} A is a \inserted{$\underline{W}$-}test
category.  In particular,  if A is totally $\underline{W}$-aspheric and i is
$M_a$-\inserted[typed]{$\underline{W}$-aspheric}\xout{\unsure{aspheric}},  A is a \inserted[typed]{strict $\underline{W}$-}\xout{\illegible{1}} test category.

\uline{Proof}   Part a) is a tautology in terms of section 51.  For part b),
to prove that $i^*(\underline{I})$ is separating,  we only have to check that for any
\xout{in $M_c$}  a in A ,  the two compositions
\diagram{\begin{tikzcd} a \arrow[r] & i^*(e_M) \simeq e_{A^\wedge} \arrow[r, shift left, "\hand{$i^*(\delta_0),i^*(\delta_1)$}"] \arrow[r, shift right] & i^*(I) \end{tikzcd}}
are distinct,  or what amounts to the same by the definition of $i^*$ ,
that the compositions
\diagram{\begin{tikzcd} i_{\hand{!}}(a) \overset{\mathrm{def}}{=} x \arrow[r] & e_M \arrow[r, shift left] \arrow[r, shift right] & I \end{tikzcd}}
are \inserted[typed]{distinct}\xout{equal}.  As  x is in $M_c$ it has a section over $e_M$ ,  so it is enough
to check that the compositions with  $e_M \longrightarrow x$  are distinct,  which just
means that  $\hand{$\delta$}_0 \hand{$\neq$} \hand{$\delta$}_1$  , qed.

\uline{Remarks}\hand{.} \hand{1)} Part b) of the corollary replaces cor.3  on\hand{/}page 253,  which is
a little monster of incongruity (as I just discovered)  -  namely two
of the assumptions on $\underline{I}$  made there (namely that $\underline{I}$\xout{r} be a multiplicative
interval,  and $I \in \mathrm{Ob}\, C$) are useless if we assume just I contractible,
moreover the \xout{assumptio} awkward separation assumption made there just
reduces by the \inserted{trivial} argument above to the assumption $\hand{$\delta$}_0 \hand{$\neq$} \hand{$\delta$}_1$ !

\hand{2)} It should be noted that the \xout{contractibilityxx} homotopy structure
on $A^\wedge$ envisioned in th.1 \unsure{of} section 79 (p.252)  is \uline{not} \deleted{thex} defined
as in the present section,  in terms of intervals in $A^\wedge$ coming from A
(call this structure h),  \uline{but as} $h' = h_{\underline{W}_A}$ defined in terms of inter-% 
\margin{formule à la ligne}
vals in $A^\wedge_{\mathrm{locas}}$\deleted{\illegible{1}} ; \inserted{this} \deleted{which} depends a priori on the choice of $\underline{W}$ ,  as it
has to because th.1  gives a criterion for the functor i to be $M_a$-$\underline{W}$-%
aspheric which does depend on $\underline{W}$ . (\over{i}{I}t surely wont be the same if we take
$\underline{W} = \underline{W}_\infty$  =  usual \inserted{weak} equivalence,  or $\underline{W} = Fl(\mathrm{Cat})$  hence $\underline{W}_A = Fl(A)$  and

\margin{30}
